\section{Problem Formulation}
\label{sec:problem}

\paragraph{Setting.}
A robot observes images $o_t$ and proprioceptive states $q_t$.
A frozen visual encoder $\varphi$ maps each image to a grid of patch features, and a context $C_t=E(h_t)$ summarizes the observed history $h_t$ using only present and past information.
At a decision time $t$, a candidate action chunk $\bA=(a_t,\dots,a_{t+H-1})$ is executed for its first $\He$ steps.
The \emph{future segment} of $\bA$ is $\tau=\big(\varphi(o_{t+1}),\dots,\varphi(o_{t+\He})\big)$, the frozen features of the observations that executing $\bA$ produces, at the native temporal resolution.
A world model must say something about $\tau$ \emph{before} $\bA$ is executed.

\paragraph{Queries.}
What a planner needs to know about $\tau$ is expressed as a family $\mathcal{Q}$ of queries.
A query $q$ has a type and parameters, such as a goal image, one or two anchor images, or a time window.
Its answer $y_q$ is computed from frozen features of the actual future, so no annotation is required.
We use five types:
\begin{enumerate*}[label=(\roman*)]
\item the similarity of the segment's end to a goal image;
\item the maximum similarity to an anchor image over the segment;
\item the number of steps from the end of the segment to a goal frame later in the same episode;
\item whether anchor $u$ is matched before anchor $v$ within the segment; and
\item masked local probes of spatial features, which retain fine motion.
\end{enumerate*}
Types (ii) and (iv) depend on the path and not only on the endpoint.
Type (iii) measures temporal progress rather than visual similarity, which matters in multi-stage tasks where the distance to a goal image is a poor measure of how far along the task is.

\paragraph{Conditional trajectory abstraction.}
A \emph{conditional trajectory code} is a sequence of $M$ discrete tokens $S\in\{1,\dots,V\}^M$ produced by a source encoder $q_\phi(S\mid C,\tau)$, which sees the actual future and is used only during training.
A reader $D_\psi(C,S,q)$ answers queries from the context and the code, and a world model $p_\theta(S\mid C,\bA)$ predicts the code from the context and the candidate chunk.
The reader never receives $\bA$, so any information about the action's consequences must pass through $S$.
In its idealized form, the code solves a conditional rate--distortion problem,
\begin{equation}
  \min_{q_\phi,\,D_\psi}\ \E_{q\sim\mathcal{Q}}\,\E\big[d\big(D_\psi(C,S,q),\,y_q\big)\big]\quad\text{s.t.}\quad I(\tau;S\mid C)\le R,
  \label{eq:crd}
\end{equation}
in which the context is available to both the encoder and the reader.
We count rate by the code itself: $M\in\{8,16,32\}$ tokens with vocabulary $V{=}256$ correspond to at most $64$, $128$, and $256$ bits per segment.
This count excludes the codebook, the reader, and the context cache, which we report separately.

\paragraph{Predictability.}
At deployment the code is not observed but predicted, $\hat S\sim p_\theta(\cdot\mid C,\bA)$, so what matters is the \emph{deployment} distortion $\E\,d\big(D_\psi(C,\hat S,q),y_q\big)$.
A code that minimizes~\eqref{eq:crd} alone may spend its rate on aspects of $\tau$ that the action does not determine, and those aspects cannot be forecast.
We therefore also require the code to be predictable, $\E[-\log p_\theta(S\mid C,\bA)]$ small, and design the code and the world model jointly (\cref{sec:codesign}).

\paragraph{Two limiting cases.}
The formulation contains the two common designs as extremes.
A \emph{frame-level world model} sets $S=\tau$: it has no codec distortion and unbounded rate, and it must forecast $\He$ frames of patch features per candidate.
A \emph{direct scorer} $D(C,\bA,q)$ removes $S$ altogether: it has zero rate, but prediction and reading are fused, so each candidate--query pair requires a separate evaluation and generalization to new queries must be learned jointly with the effect of the action.
\method operates at a finite, explicit rate between them.
The questions we study are whether, at such a rate, predicted codes answer queries as accurately as a frame-level model at lower cost, and whether they generalize to held-out goals better than a direct scorer.
\Cref{tab:designs} summarizes the design properties.

\paragraph{Policy-guided planning.}
We evaluate world models in the planning interface of recent policy steering work~\cite{qi2025gpc,cui2026dreamsteer}.
At every decision, a frozen pretrained policy $\pi$ proposes $K$ chunks $\bA^{(1)},\dots,\bA^{(K)}$ from independent noise seeds, where $\bA^{(1)}$ is the chunk the unmodified policy would execute.
Each world model scores every proposal with the same goal query, the highest-scoring proposal is executed for $\He$ steps with ties resolved in favor of $\bA^{(1)}$, and the policy replans.
Selecting among proposals of a competent policy, rather than optimizing actions against the model, keeps candidates on the policy's support, where a learned model is less exposed to adversarial exploitation of its errors~\cite{janner2019mbpo,gao2023overoptimization}.
The same proposal bank is given to every compared model.
Because no world model can recover an action that is absent from the bank, our evaluation first verifies that the bank contains better candidates (\cref{sec:experiments}).

\begin{table}[t]
\centering
\small
\setlength{\tabcolsep}{3pt}
\resizebox{\linewidth}{!}{%
\begin{tabular}{@{}lcccc@{}}
\toprule
 & Predicted per & Sequential & New goal & Path \\
 & candidate & steps & or query & queries \\
\midrule
Frame-level WM & $\He\cdot P$ tokens & $\He$ & reader only & \checkmark \\
Endpoint latent & $P$ tokens & $1$ & reader only & -- \\
Direct scorer & one score & $1$ per query & new evaluation & trained only \\
\method (ours) & $M$ tokens & $M$ & reader only & \checkmark \\
\bottomrule
\end{tabular}}
\caption{\textbf{Design properties} (not results) of world models for evaluating a candidate chunk. $P$ is the number of patch tokens per frame (hundreds), $M\in\{8,16,32\}$ the number of code tokens. \method's $M$ autoregressive steps are short but not automatically cheaper than $\He$ frame steps, so latency is measured, not assumed.}
\label{tab:designs}
\end{table}
